\chapter{Wprowadzenie}
\label{cha:wstep}

\textbf{REPAIRO} to aplikacja desktopowa, zaprojektowana do zarządzania serwisem samochodowym. Powstała w języku Java z wykorzystaniem biblioteki JavaFX oraz relacyjnej bazy danych MySQL. System umożliwia obsługę procesów związanych z funkcjonowaniem warsztatu: rejestrację klientów i pojazdów, tworzenie zleceń napraw, zarządzanie usługami i częściami oraz przypisywanie pracowników do konkretnych zadań.

Interfejs użytkownika został zaprojektowany w duchu przejrzystości i intuicyjności - wszystkie funkcjonalności dostępne są z poziomu czytelnych formularzy, a dane wprowadzane przez użytkownika są walidowane, co ogranicza ryzyko błędów. System automatycznie przelicza koszt zlecenia na podstawie wybranych usług i części, wspierając dokładność wycen oraz organizację pracy.

Aplikacja może być wykorzystywana zarówno w małych, jak i średnich warsztatach, gdzie potrzeba systematycznego zarządzania danymi klientów i pojazdów staje się coraz bardziej istotna.\\

\textbf{REPAIRO} is a desktop application designed to manage a car service. It was created in Java using the JavaFX library and the MySQL relational database. The system enables the handling of processes related to the functioning of the workshop: registration of customers and vehicles, creation of repair orders, management of services and parts, and assignment of employees to specific tasks.

The user interface was designed in the spirit of transparency and intuitiveness - all functionalities are available from the level of legible forms, and data entered by the user is validated, which reduces the risk of errors. The system automatically calculates the cost of the order based on selected services and parts, supporting the accuracy of pricing and work organization.

The application can be used in both small and medium-sized workshops, where the need for systematic management of customer and vehicle data is becoming increasingly important.